\subsection{Local probabilities and the conditional GP global studies}
\label{sec:v5-global-results}

The following plots apply the method in Section~\ref{sec:v5-global} to each
campaign's complete mass scan, including 2015 through 100~MeV. Their top panels show the observed signed fit
$r$, the unfluctuated stress response $a_m$, and the mean of complete Poisson scans. The updated 2015 and combined
displays use the same 256 coherent experiments at every mass; the unchanged
2016 and 2021 displays retain their 1,000-scan benchmark. The lower panels put probabilities and their one-sided normal
quantiles on the same mass axis. Black denotes the conditional local result
and red the conditional GP global result; gray retains the nominal local
asymptotic comparison. The global probability comes from the maximum of a
correlated field with the declared positive-fit gate. It is different from
the Sidak reference accompanying the asymptotic plots.


The implementation comparison, stress-reference construction, and treatment of
rare Gaussian tails are documented in Appendix~\ref{app:v504-global-study}.
Figures~\ref{fig:v502-stress-distributions} and
\ref{fig:v502-injection-recovery} there show why large reference offsets can
coexist with ordinary observed fits and why injected signals do not remove
those offsets. The extended 2015 and combined fields are recomputed with the
same search membership as the observed limits, rather than assigning the old
global penalty to the enlarged search.

\paragraph{The apparent 76 MeV spike and the detector resolution.}
The sharp excursion belongs to the \emph{combined, stress-centered} display.
The 2016 raw fit at 76~MeV is negative ($r=-2.46$), while the combined raw
response is only $r=0.166$. On the saved 1 MeV test-mass grid, the combined signed fit changes sign
between 75 and 76~MeV, opening the positive-fit gate while the stress mean is strongly
negative. The resulting spike in probability or $Z$ is therefore not a
measured resonance width. The adopted 2016 resolution at 76~MeV is
$\sigma_m=3.10$~MeV (Gaussian FWHM 7.29~MeV). The positive raw response near
90--92~MeV warrants examination of the mass residual and its background
dependence; neither its visual plausibility nor its stress-centered ranking
establishes a physical signal. Appendix~\ref{app:v504-width-review} gives
the signed-root decomposition and the saved width and injection controls.

\paragraph{Dependence on binning.}
The zero crossing refers to a scan of fitted amplitudes, not to a zero in
the measured histogram. Changing histogram bins changes the GP training
counts, the excluded-bin membership and the integrated signal template;
changing the test-mass grid only changes where that response is sampled.
The v5.1.0 fixed-kernel rebinning study gives 2016 signed roots at 76 MeV of
$-2.460$, $-2.545$ and $-2.432$ for 0.25, 0.5 and 1 MeV bins. Its stress
responses are $-14.653$, $-14.620$ and $-14.488$. Thus coarsening those bins
does not remove the 2016 deficit or the large stress offset. The corresponding
2021 roots are $1.808$, $1.429$ and $1.517$ for 0.625, 1.25 and 2.5 MeV bins.

A bounded joint check in v5.0.5 repeats the signed shared-coupling fit at
72--80 MeV using the same spectra, detector widths, native-density conversion
and frozen kernel coordinates. With the bin origins unchanged, doubling
and quadrupling the histogram widths changes the combined root at 76 MeV
from $+0.166$ to $-0.216$ and $-0.137$. Shifting each coarse origin by half
its new bin width gives $+0.530$ and $-0.695$, respectively. The first
positive tested point is consequently 76 or 77 MeV in these alternatives.
The zero crossing is sensitive to binning at the one-step level.

Figure~\ref{fig:v505-joint-binning} and
Appendix~\ref{app:v505-joint-binning} give the full comparison, including
the small accompanying changes in sampled support endpoints. These checks
recompute the observed pointwise fits; they do not recompute the stress
field or its global tail. The spike's position and height therefore cannot
be interpreted as a detector-resolved resonance feature. Binning should be
chosen from resolution and closure controls, before inspecting which choice
makes the observed probability more striking.

\paragraph{What the stress response means.}
The gold curve is the signed fit obtained by analyzing a smooth expected
background with no injected signal and no Poisson fluctuation. It measures
how that particular background shape is represented by the moving GP
interpolation and signal fit. It is not an observed significance and is not
assumed to equal the exact Poisson mean; the green curve checks the latter.
A response near zero is desirable for an unbiased signal search. A large
response can occur when the chosen background leaves a signal-shaped excess
or deficit relative to the interpolation. More simulations can measure that
response more precisely, but cannot make the generating shape appropriate
for the data.

\fig{0.98}{../figures/v5_global_probabilities_2015.pdf}{Profiled 2015 conditional local and global probabilities on 19--100~MeV.
The full field uses 82 integer hypotheses and the endpoint kernel prescription
in Appendix~\ref{app:v504-extension}. Its response matrix and the first 256
coherent Poisson scans are reused from the extended covariance study;
200,000 Gaussian fields integrate the enlarged search. Open symbols retain
one-sided 95\% bounds for empty direct tails. Probability and $Z^+$ panels
refer to the declared stress model. The interpretation, sigcorr check and
sampling method are in Appendix~\ref{app:v504-global-study}.}
{fig:v5-global-2015}
\fig{0.98}{../figures/v5_global_probabilities_2016.pdf}{Profiled 2016 conditional
global study with the same black-local/red-global convention and integer
$Z^+$ guides. Ordinary sampling uses 2,000,000 fields; targeted importance
sampling resolves the previously empty Gaussian tails without altering the
field model or ordering. At 42~MeV the conditional Gaussian global
probability is $6.096\times10^{-20}$ ($Z=9.067$), with 0.16\% relative Monte
Carlo standard error. This extreme value reflects the biased stress
reference, not a particle discovery. The direct Poisson check remains
0/1,000, giving only $p<0.002991$ ($Z>2.749$); its open triangles retain
those bounds. Red shading gives sampling intervals, not background-model
uncertainty. This global calculation is different from the Sidak reference. The implementation and stress-reference diagnostics are in Appendix~\ref{app:v504-global-study}.}
{fig:v5-global-2016}
\fig{0.98}{../figures/v5_global_probabilities_2021.pdf}{Profiled study for the
released 2021 10\% sample. The principal ordering selects 92~MeV and gives
GP global probability 0.02519 ($Z=1.957$), with 26/1,000 direct exceedances.
The GP ensemble retains 200,000 fields. Black is the conditional local
probability and red the global probability, with the corresponding $Z^+$
below. This global calculation uses the correlated field rather than the
Sidak correction used for the asymptotic references. Qualification of the
generating background remains required for a particle-global interpretation. The implementation and stress-reference diagnostics are in Appendix~\ref{app:v504-global-study}.}
{fig:v5-global-2021}
\FloatBarrier

\fig{0.98}{../figures/v5_global_probabilities_combined.pdf}{Combined conditional field over the complete 19--250~MeV search,
with all three campaigns active through 100~MeV. The ten changed coordinates
use new joint-likelihood responses to the same 1,626 input-bin directions;
all 256 direct experiments remain coherent over the entire grid. Gaussian
tails use 200,000 ordinary fields and 40,000 conditional-mixture draws at
every positive-fit threshold with standardized score at least 3.5. At
76~MeV the conditional global probability is $1.351\times10^{-17}$, but the
raw fit is only $r=0.166$: this is a stress-reference mismatch. Empty direct
tails give only $p<0.01163$. Appendix~\ref{app:v504-global-study} explains
the estimator, and Appendix~\ref{app:v504-width-review} explains the narrow
display feature. Historical values for the former membership are retained
in Appendix~\ref{app:v504-historical-displays}.}
{fig:v5-global-combined}
\FloatBarrier

For comparison with the conditional tests in Appendix~\ref{app:v504-global-study},
Figure~\ref{fig:v502-combination-discovery} completes the nominal discovery
reference plots for every pair and the all-three overlap. It keeps each
scope's local fit and its own Sidak search penalty together.

\fig{0.99}{../figures/v502_combination_discovery.pdf}{Discovery probabilities
and aligned $Z^+$ values for all four overlapping combinations. Local
asymptotic curves are black and Sidak global references are red. Each pair
of panels shares its mass axis and uses one fixed $N_{\rm eff}$ over that
scope's full overlap, now through 100~MeV whenever 2015 participates.
The 90--100~MeV extension is shaded. These are analytic Sidak references, distinct from the stress-centered
correlated-field global probabilities above; no additional correction for
selecting among overlapping scopes is included.}
{fig:v502-combination-discovery}
\FloatBarrier

\paragraph{Fixed-mass recovery and full-scan signal responses.}
The injection masses are inherited diagnostic anchors: 51, 78 and 95~MeV
for 2015; 42, 66, 76 and 92~MeV for 2016; and 71 and 92~MeV for 2021.
They cover the earlier 2016 stress excursions, the neighboring structures
illustrated by the 2021 65/71/78~MeV examples, representative peaks, and
the 2015 extension. They were chosen to investigate those known
features, so they do not constitute an independent discovery selection.
Each mass is tested against both the archived coherent stress spectrum
and a full-support local GP control fitted to the observed sidebands
with that signal window excluded. The control depends on the injection
mass, so these local controls do not form a single global null. The same frozen reference
error defines the physical yields $A/\sigma_{\rm ref}=0,2,5$ for both
truths. Within a mass and truth, 64 Poisson background spectra are paired
across strengths; independent Poisson signal increments of two and three
reference errors build the size-two and size-five samples. Thus the
three strengths are correlated experiments, not independent ensembles.
Figure~\ref{fig:v511-injections} in Appendix~\ref{app:v504-global-study}
shows three representative anchors, including the 2015 95~MeV extension check.
At 2016 76~MeV, the size-five injection gives mean signed root $-9.80$
under the stress truth and $+4.52$ under the local GP control; every
stress fit remains negative. At 2021 71~MeV the corresponding means are
$+3.14$ and $+4.37$. Adding a signal moves the distributions, but it does
not remove the background-dependent offset. With 64 experiments per
curve, these are small conditional recovery samples.

A deterministic injected spectrum is also scanned over every tested mass.
For a count perturbation $\delta n_i$ from the signal, the first-order
prediction from the background response is
\begin{equation}
 \Delta r_m^{\rm lin}=\sum_{i:B_i>0}D_{im}
                         \frac{\delta n_i}{\sqrt{B_i}}.
 \label{eq:v511-injection-projection}
\end{equation}
For the 2016 76~MeV size-five injection, the exact change is $+4.843$
at the injected mass and reaches $-3.589$ at 85~MeV. The largest difference
from the linear projection is 0.00924 in signed root over the scan.
Comparison with the exact difference of the injected and background-only
scans tests linearity and reveals neighboring peaks or deficits induced by
the shared interpolation. An off-mass dip can therefore accompany a
positive injection. This diagnostic does not establish the origin of a
feature in the observed spectrum. A covariance-matrix column describes
shared noise; it is not the signal-template projection in
Eq.~\ref{eq:v511-injection-projection}.
\fig{0.99}{../figures/v511_signed_injection_echoes.pdf}{Deterministic signal
injections into the coherent stress spectrum. Left: signed roots before
and after injection. Right: the size-five change with full retraining,
its first-order template projection, and a control holding the
background-only GP prediction and covariance fixed. The likelihood's
background nuisance parameters remain profiled in the last control.
Vertical lines mark the injected mass. Neighboring negative responses
are retained. These full scans contain no Poisson fluctuation.}
{fig:v511-echo}
\FloatBarrier



\FloatBarrier\clearpage

\subsection{Agreement of fitted rates across campaigns}
\label{sec:v5-rate-consistency}

A shared coupling requires the signal yields in different campaigns to agree
after their different normalizations are applied. Figure~\ref{fig:v5-consistency}
compares the separate fitted coupling estimates with the common fit. The
quantity formerly labelled $\Delta D$ is the loss in twice the log likelihood
when separate campaign amplitudes are replaced by one shared amplitude:
\begin{equation}
 \Delta D=2\left[\mathrm{NLL}_{\rm common}-
                    \sum_d\mathrm{NLL}_{d,\rm independent}\right].
\end{equation}
A larger loss means the common rate describes the selected local spectra less
well than separate rates. The number of freed amplitudes is one fewer than
the number of active campaigns. Because the displayed masses were selected
from the data, this descriptive comparison is not assigned an unadjusted
compatibility probability.

\widefig{0.98}{../figures/v5_dataset_amplitude_consistency.pdf}{Separate campaign coupling estimates with local curvature errors and the
common fitted rate at 66, 92, 72 and 76~MeV. All three campaigns now
enter every comparison. At 92~MeV the common fit gives
$\widehat\epsilon^2=3.096\times10^{-6}$; freeing the three rates lowers
twice the negative log likelihood by 11.552, compared with 2.845 at 66~MeV.
The shaded band is the common fit's local curvature interval. These selected
comparisons are descriptive; the errors are not post-selection intervals and
no unadjusted compatibility probability is assigned.}
{fig:v5-consistency}
\FloatBarrier


The extension makes the 92~MeV comparison especially informative. In units
of $10^{-6}$, the separate fitted rates and local curvature errors are
$77.48\pm29.96$ for 2015, $9.13\pm2.93$ for 2016, and
$1.65\pm1.35$ for 2021. The 2015 rate is weakly constrained, but its central
value is much larger than the common estimate $3.096\pm1.229$.
The shared fit therefore does not extract each campaign's preferred peak
height independently. Its nominal local $Z=2.520$ and 90\% \CLs{} upper
limit $4.675\times10^{-6}$ summarize a common-rate model, with its
normalizations and background constraints, rather than three separately
confirmed excesses. The likelihood loss is 11.552 for two additional free
amplitudes; at 66~MeV the analogous loss is 2.845. This comparison identifies
rate agreement as a useful diagnostic without assigning a probability to
masses chosen from the same scans.

The 2015 template width rises from 4.70~MeV at 90~MeV to 5.23~MeV at
100~MeV under the inherited linear response. Broader resolution makes nearby
hypotheses strongly dependent, but it does not explain away a rate mismatch
or establish that a visible structure is physical. A credible common signal
would need stable rates after independently supported background and resolution
variations, a signal shape consistent with the detector response, and a
predeclared global treatment of the selected masses and model choices.
The conditional GP probabilities, nominal asymptotic $p_0$, upper-limit tail
diagnostics, and common-rate comparison answer different questions; none
provides a probability that the 90--92~MeV feature is a new particle.

\clearpage
